\backhead{FURTHER EXERCISES}{khred}
\begin{exlist}
\item Find the drain spacing by Hooghoudt for $K=1.2$ m\,d$^{-1}$, $d_e=2.0$ m, $h=0.6$ m and $q=10$ mm\,d$^{-1}$.\exans{36.4 m}
\item Repeat the previous exercise with $d_e=1.0$ m and comment on the effect of the depth of the impermeable layer.\exans{27.4 m}
\item With the spacing of the first exercise, $\mu=0.06$ and $D=2.3$ m, find the reaction factor of Glover--Dumm and the time for the water table to fall to a quarter of its height.\exans{0.343 d$^{-1}$; 4.5 d}
\item A drain of 100 mm diameter at a slope of 0.002 must carry the discharge of a field of 300 m by 40 m at 10 mm\,d$^{-1}$. Find the discharge in L\,s$^{-1}$.\exans{1.39 L\,s$^{-1}$}
\item Irrigation water of salinity 1.2 dS\,m$^{-1}$ is applied to a crop tolerating 4 dS\,m$^{-1}$ in the drainage water. Find the leaching requirement $LR=EC_i/EC_d$ and the gross depth for a net 500 mm.\exans{0.30; 714 mm}
\end{exlist}
\backhead{ESSAY QUESTIONS}{khred}
{\small\itshape Answer each question in about one page, with equations, a sketch and at least one number from this chapter.}\par\vspace{3pt}
\begin{exlist}
\item Explain the Hooghoudt equation and the meaning of the equivalent depth.
\item Discuss the steady and the unsteady approaches to drain design, and when each is appropriate.
\item Explain why drainage and irrigation must be designed together in arid regions.
\item Discuss controlled drainage and its benefits for water and nitrogen.
\item Explain the leaching requirement and its link to the salinity of the irrigation water.
\item Discuss how the same drains behave in the ten soils of the study.
\end{exlist}
